\clearpage\pagestyle{otMain}\setlength{\otBlocEspacement}{2.5mm}\fontsize{11}{15}\selectfont\corrigesChapitre{11}{Probabilités : événements et deux épreuves}
\begin{blocLivre}{cor-11-1}\begin{corrige}{1}{Décrire un univers}\demarche{Écrire les issues avant de compter.}
L’univers est $\res{\{1\,;2\,;3\,;4\,;5\,;6\,;7\,;8\}}$. Les multiples de 3 sont 3 et 6 : $A=\res{\{3\,;6\}}$. Les huit issues sont équiprobables, donc $P(A)=\frac28=\res{\frac14}$.
\end{corrige}\end{blocLivre}
\begin{blocLivre}{cor-11-2}\begin{corrige}{2}{L’événement contraire}\demarche{Prendre le complémentaire dans l’univers.}
$\overline A$ contient les issues qui ne sont pas dans $A$ : $\res{\{1\,;2\,;4\,;5\,;7\,;8\}}$. Par comptage, $P(\overline A)=\frac68=\frac34$. Par le contraire, $1-\frac14=\res{\frac34}$.
\end{corrige}\end{blocLivre}
\begin{blocLivre}{cor-11-3}\begin{corrige}{3}{Certain et impossible}\demarche{Comparer l’événement à l’univers.}
« Obtenir un nombre inférieur à 7 » contient toutes les issues : il est certain, de probabilité $\res1$. « Obtenir 8 » ne contient aucune issue : c’est $\varnothing$, l’événement impossible, de probabilité $\res0$.
\end{corrige}\end{blocLivre}
\begin{blocLivre}{cor-11-4}\begin{corrige}{4}{Plusieurs écritures}\demarche{Passer de la fraction au quotient puis au pourcentage.}
$\frac38=\res{0{,}375}$, soit $\res{37{,}5\,\%}$. Les trois écritures représentent la même probabilité.
\end{corrige}\end{blocLivre}
\begin{blocLivre}{cor-11-5}\begin{corrige}{5}{Une urne}\demarche{Compter les objets équiprobables.}
Les 8 boules ont la même chance d’être tirées. $P(R)=\frac48=\res{\frac12}$. $\overline R$ contient les 3 bleues et la verte : $P(\overline R)=\frac48=\res{\frac12}$. Compter trois couleurs ne donnerait pas trois chances égales.
\label{fig-cor-11-5}\figureGeo[corrige]{urne-431}
\end{corrige}\end{blocLivre}
\begin{blocLivre}{cor-11-6}\begin{corrige}{6}{Intersection et réunion}\demarche{Placer chaque issue dans le bon ensemble.}
$A=\{2\,;4\,;6\}$ et $B=\{4\,;5\,;6\}$, de probabilité $\frac12$ chacun. $A\cap B=\res{\{4\,;6\}}$ : probabilité $\res{\frac13}$. $A\cup B=\res{\{2\,;4\,;5\,;6\}}$ : probabilité $\frac46=\res{\frac23}$. Les issues 4 et 6, communes, ne sont comptées qu’une fois dans la réunion.
\label{fig-cor-11-6}\figureGeo[corrige]{venn-de}
\end{corrige}\end{blocLivre}
\begin{blocLivre}{cor-11-7}\begin{corrige}{7}{Deux pièces}\demarche{Distinguer les deux épreuves.}
L’arbre donne quatre issues équiprobables : $(P\,;P)$, $(P\,;F)$, $(F\,;P)$ et $(F\,;F)$. L’événement contient $(P\,;F)$ et $(F\,;P)$ : sa probabilité vaut $\frac24=\res{\frac12}$. Les deux pièces sont distinguées, donc $PF$ et $FP$ sont deux issues différentes.
\label{fig-cor-11-7}\figureGeo[corrige]{arbre-deux-pieces}
\end{corrige}\end{blocLivre}
\begin{blocLivre}{cor-11-8}\begin{corrige}{8}{Une pièce et un dé}\demarche{Repérer les cases du tableau qui conviennent.}
Le tableau compte $2\times6=\res{12}$ issues équiprobables. « Pile et pair » contient $(P\,;2)$, $(P\,;4)$ et $(P\,;6)$ : probabilité $\frac3{12}=\res{\frac14}$. « Face ou 6 » contient les six issues avec face et $(P\,;6)$ : probabilité $\res{\frac7{12}}$.
\label{fig-cor-11-8}\figureGeo[corrige]{tableau-piece-de}
\end{corrige}\end{blocLivre}
\begin{blocLivre}{cor-11-9}\begin{corrige}{9}{La somme de deux dés}\demarche{Compter les cases du tableau à double entrée.}
Le tableau contient 36 issues équiprobables. La somme 7 apparaît 6 fois : probabilité $\frac6{36}=\res{\frac16}$. Les sommes 10, 11 et 12 apparaissent $3+2+1=6$ fois : probabilité $\res{\frac16}$ aussi.
\label{fig-cor-11-9}\figureGeo[corrige]{somme-deux-des}
\end{corrige}\end{blocLivre}
\begin{blocLivre}{cor-11-10}\begin{corrige}{10}{Des issues non équiprobables}\demarche{Utiliser la somme des probabilités des issues.}
Le total des trois probabilités vaut 1, donc $p=1-0{,}2-0{,}5=\res{0{,}3}$. La probabilité de ne pas obtenir B est $1-0{,}5=\res{0{,}5}$. On ne compte pas simplement un tiers par issue.
\end{corrige}\end{blocLivre}
\begin{blocLivre}{cor-11-11}\begin{corrige}{11}{Une observation}\demarche{Séparer observation et modèle.}
La fréquence observée est $\frac{31}{50}=\res{0{,}62}$. Le modèle reste une probabilité de pile égale à $\res{\frac12}$ pour le prochain lancer. Sur un nombre fini de lancers, la fréquence peut s’écarter de la probabilité.
\end{corrige}\end{blocLivre}
\begin{blocLivre}{cor-11-12}\begin{corrige}{12}{Quatre séries de lancers}\demarche{Diviser chaque effectif par le nombre de lancers.}
Les fréquences sont $\res{0{,}46}$, $\res{0{,}56}$, $\res{0{,}50}$ et $\res{0{,}62}$. Pour un même nombre de lancers, le hasard fait varier les résultats : c’est la fluctuation. Sur les 200 lancers, $\frac{107}{200}=\res{0{,}535}$, proche de la probabilité 0,5.
\label{fig-cor-11-12}\figureGeo[corrige]{fluctuation-pieces}
\end{corrige}\end{blocLivre}
\begin{blocLivre}{cor-11-13}\begin{corrige}{13}{Fréquences et probabilités}\demarche{Associer chaque nombre de piles à ses issues.}
Fréquences : $\frac{13}{60}\approx\res{0{,}22}$, $\frac{32}{60}\approx\res{0{,}53}$ et $\frac{15}{60}=\res{0{,}25}$. Le modèle des quatre issues équiprobables donne $\frac14$, $\frac12$ et $\frac14$. Les fréquences en sont proches sans leur être égales : les deux diagrammes ont la même allure.
\label{fig-cor-11-13}\figureGeo[corrige]{distribution-deux-pieces}
\end{corrige}\end{blocLivre}
\begin{blocLivre}{cor-11-14}\begin{corrige}{14}{Un effectif inconnu}\demarche{Retrouver le total à partir de la part rouge.}
Les rouges représentent un quart du total, donc le sac contient $5\times4=20$ boules. Les bleues sont $20-5=\res{15}$. Contrôle : $\frac5{5+15}=\frac14$.
\end{corrige}\end{blocLivre}
\begin{blocLivre}{cor-11-15}\begin{corrige}{15}{Des numéros et des couleurs}\demarche{Énumérer les issues favorables sans répétition.}
$R\cap U=\res{\{R1\}}$ et $R\cup U=\res{\{R1\,;R2\,;R3\,;B1\,;V1\}}$, donc $P(R\cup U)=\res{\frac56}$. Or $P(R)+P(U)=\frac36+\frac36=1$ : le jeton R1, qui appartient aux deux événements, serait compté deux fois.
\label{fig-cor-11-15}\figureGeo[corrige]{jetons-numerotes}
\end{corrige}\end{blocLivre}
